\documentclass[11pt]{article}
\usepackage[margin=1in]{geometry}
\usepackage{amsmath,amssymb,amsthm}
\usepackage[hidelinks]{hyperref}

\newtheorem{theorem}{Theorem}

\title{A Partial Answer to Baudier's \(C(K)\)-Universality Question}
\author{agent\_lane\_06}
\date{2026-06-29}

\begin{document}
\maketitle

\section*{Status}

\textbf{partial\_result\_likely\_valid.}  This packet answers the isometric-universality clause of Question 2 in arXiv:1305.4025 for compact metric spaces, and gives the uncountable compact metric side of the \(c_{C(K)}=1\) question.  It does not settle the remaining almost-isometric \(c_{C(K)}=1\) problem for countable compact metric spaces of high Cantor-Bendixson rank.

\section*{Source Question}

The source is Florent Baudier, \emph{A topological obstruction for small-distortion embeddability into spaces of continuous functions on countable compact metric spaces}, arXiv:1305.4025.  In Section 3, Question 2 asks:
\[
\text{Characterize compacta }K\text{ such that }c_{C(K)}(\mathrm{SEP})=1
\]
or \(c_{C(K)}(\mathcal{SEP})=1\), and asks for which \(K\) the space \(C(K)\) is an isometric universal Banach space for the class of separable Banach spaces.

\section*{Idea of the Proof}

The two sides are separated by the perfect-set theorem.  An uncountable compact metric space contains a closed Cantor set.  The Cantor set maps continuously onto \([0,1]\), so \(C([0,1])\) embeds isometrically into \(C\) of that Cantor set by pullback.  The Borsuk-Dugundji extension theorem then embeds that \(C\)-space isometrically into \(C(K)\).  Thus the Banach-Mazur theorem transfers the usual universality of \(C([0,1])\) to \(C(K)\).

For countable compact metric \(K\), the dual \(C(K)^*\) is separable, since measures on the scattered countable compactum are atomic and form \(\ell_1(K)\).  Hence no subspace of \(C(K)\) can be linearly isomorphic to \(\ell_1\).  Godefroy-Kalton's isometric lifting theorem converts a metric isometric copy of \(\ell_1\) in any Banach space into a linear isometric copy, so countable compacta are excluded.

\begin{theorem}
Let \(K\) be a compact metric space.  Then \(C(K)\) is an isometric universal Banach space for the class of separable Banach spaces if and only if \(K\) is uncountable.

In addition, if \(K\) is uncountable compact metric, then
\[
c_{C(K)}(\mathrm{SEP})=c_{C(K)}(\mathcal{SEP})=1.
\]
\end{theorem}

\begin{proof}
Assume first that \(K\) is uncountable compact metric.  By the classical perfect-set theorem, \(K\) contains a closed subset \(F\) homeomorphic to the Cantor set \(2^{\mathbb N}\).  There is a continuous surjection \(q:F\to[0,1]\).  Therefore
\[
J:C([0,1])\to C(F),\qquad Jf=f\circ q,
\]
is a linear isometry, because \(q(F)=[0,1]\).

By the Borsuk-Dugundji extension theorem, since \(F\) is closed in the compact metric space \(K\), there is a linear extension operator
\[
E:C(F)\to C(K)
\]
with \(\|E\|=1\) and \((Eg)|_F=g\) for every \(g\in C(F)\).  It follows that \(E\) is an isometry: \(\|Eg\|_{C(K)}\ge \|g\|_{C(F)}\) by restriction to \(F\), while \(\|Eg\|\le\|g\|\) by \(\|E\|=1\).  Hence \(EJ\) is a linear isometric embedding of \(C([0,1])\) into \(C(K)\).

The Banach-Mazur theorem says that every separable Banach space embeds linearly isometrically into \(C([0,1])\), so every separable Banach space embeds linearly isometrically into \(C(K)\).  Thus \(C(K)\) is isometrically universal for separable Banach spaces.  Also, every separable metric space embeds isometrically into a separable Lipschitz-free space, and this Banach space embeds linearly isometrically into \(C([0,1])\).  Hence every separable metric space embeds isometrically into \(C(K)\), proving \(c_{C(K)}(\mathrm{SEP})=c_{C(K)}(\mathcal{SEP})=1\).

Conversely, suppose \(K\) is countable compact metric.  Then \(K\) is scattered, and every regular Borel measure on \(K\) is purely atomic.  Thus \(C(K)^*\) is isometric to \(\ell_1(K)\), which is separable because \(K\) is countable.  If \(Y\subseteq C(K)\) is a closed subspace, then \(Y^*\) is a quotient of \(C(K)^*\), so \(Y^*\) is separable.  Therefore \(C(K)\) cannot contain a linear subspace isomorphic to \(\ell_1\), since \(\ell_1^*=\ell_\infty\) is nonseparable.

If \(C(K)\) were isometrically universal for separable Banach spaces, then \(\ell_1\) would embed isometrically into \(C(K)\) as a metric space.  By the theorem of Godefroy and Kalton, whenever a separable Banach space embeds isometrically into a Banach space, it is linearly isometric to a subspace of that Banach space.  This would give a linear isometric copy of \(\ell_1\) in \(C(K)\), contradicting the previous paragraph.  Hence countable compact metric \(K\) are not isometric universal Banach compacta.
\end{proof}

\section*{Limitations}

This theorem does not decide whether there are countable compact metric spaces \(K\) with \(c_{C(K)}(\mathrm{SEP})=1\) or \(c_{C(K)}(\mathcal{SEP})=1\) in the almost-isometric sense.  Baudier's Corollary 2.4 only forces \(\sigma(K)\ge\omega\) for such examples; the countable high-rank case remains outside this packet.

\section*{Novelty and Search Notes}

The lightweight run indexes had no exact entry for arXiv:1305.4025.  Local and web searches for the source title, \(c_{C(K)}(\mathrm{SEP})\), and \(C(K)\)-distortion did not reveal a later exact packet or answer in this run.  The argument uses standard classical ingredients: the perfect-set theorem for uncountable compact metric spaces, the Borsuk-Dugundji extension theorem, Banach-Mazur universality of \(C([0,1])\), and the Godefroy-Kalton theorem on isometric embeddings of separable Banach spaces.

\section*{References}

\begin{itemize}
\item F. Baudier, \emph{A topological obstruction for small-distortion embeddability into spaces of continuous functions on countable compact metric spaces}, arXiv:1305.4025.
\item S. Banach and S. Mazur, \emph{Zur Theorie der linearen Dimension}, Studia Math. 4 (1933), 100--112.
\item J. Dugundji, \emph{An extension of Tietze's theorem}, Pacific J. Math. 1 (1951), 353--367.
\item G. Godefroy and N. J. Kalton, \emph{Lipschitz-free Banach spaces}, Studia Math. 159 (2003), 121--141.
\end{itemize}

\end{document}
